\chapter{范畴语言初步}

范畴语言起初看上去像是一套抽象记号，但它的真正作用并不在于追求抽象本身，而在于把“对象、对象之间的映射、以及这些映射如何彼此兼容”统一地组织起来。在线性代数中，我们研究向量空间与线性映射；在群论中，我们研究群与群同态；在模形式理论中，我们研究模曲线上的层、上同调、Hecke 对应以及它们诱导的算子。范畴语言正是描述这些现象的共同语法。

本附录只介绍最基本的概念：范畴、函子、自然变换、正合序列与正合函子，并在最后说明它们为何在现代模形式理论中不可或缺。读者不必在第一次阅读时掌握一切细节，但应熟悉这些定义及典型例子。

\section{范畴与函子}

\subsection{范畴的定义}

\begin{definition}
一个\emph{范畴} $\mathcal C$ 由以下数据组成：
\begin{enumerate}[label=(\arabic*), leftmargin=2.5em]
  \item 一个对象类，记其元素为 $X,Y,Z,\dots$；
  \item 对任意两个对象 $X,Y$，给定一个态射集合 $\Hom_{\mathcal C}(X,Y)$，其元素称为从 $X$ 到 $Y$ 的\emph{态射}，记作 $f\colon X\to Y$；
  \item 对任意三个对象 $X,Y,Z$，给定复合运算
  \[
    \Hom_{\mathcal C}(Y,Z)\times \Hom_{\mathcal C}(X,Y)\longrightarrow \Hom_{\mathcal C}(X,Z),
    \qquad (g,f)\longmapsto g\circ f;
  \]
  \item 对任意对象 $X$，给定一个恒等态射 $\Id_X\colon X\to X$。
\end{enumerate}
这些数据满足：
\begin{enumerate}[label=(\roman*), leftmargin=2.5em]
  \item \textbf{结合律：}若 $f\colon W\to X$，$g\colon X\to Y$，$h\colon Y\to Z$，则
  \[
    h\circ (g\circ f)=(h\circ g)\circ f;
  \]
  \item \textbf{恒等律：}若 $f\colon X\to Y$，则
  \[
    f\circ \Id_X=f,
    \qquad
    \Id_Y\circ f=f.
  \]
\end{enumerate}
\end{definition}

\begin{remark}
严格地说，“对象的全体”常常不是一个集合，而是一个类；不过在本书中我们不必深究集合论基础。对初学者而言，把范畴看作“对象与保持结构的映射所构成的系统”即可。
\end{remark}

\begin{example}
集合范畴 $\mathbf{Set}$ 的对象是一切集合，态射是集合之间的映射，复合就是通常的函数复合，恒等态射就是恒等映射。
\end{example}

\begin{example}
群范畴 $\mathbf{Grp}$ 的对象是一切群，态射是群同态。类似地，环范畴 $\mathbf{Ring}$ 的对象是一切含幺环，态射是保持加法、乘法与单位元的环同态。
\end{example}

\begin{example}
给定一个域 $\F$，向量空间范畴 $\mathbf{Vect}_{\F}$ 的对象是所有 $\F$-向量空间，态射是 $\F$-线性映射。
\end{example}

\begin{example}
拓扑空间范畴 $\mathbf{Top}$ 的对象是拓扑空间，态射是连续映射。
\end{example}

\begin{example}
设 $R$ 为环。左 $R$-模范畴 $R\text{-}\mathbf{Mod}$ 的对象是左 $R$-模，态射是 $R$-模同态。这是本附录中最重要的例子之一，因为正合序列与导出函子主要在模范畴中出现。
\end{example}

\begin{definition}
设 $\mathcal C$ 为范畴。若态射 $f\colon X\to Y$ 存在态射 $g\colon Y\to X$ 使得
\[
  g\circ f=\Id_X,
  \qquad
  f\circ g=\Id_Y,
\]
则称 $f$ 是一个\emph{同构}，并称对象 $X$ 与 $Y$ 在 $\mathcal C$ 中\emph{同构}，记作 $X\cong Y$。
\end{definition}

\begin{remark}
范畴论中常常只关心对象到同构为止的性质。比如在线性代数里，有限维 $\F$-向量空间按维数完全分类；换句话说，在 $\mathbf{Vect}_{\F}$ 中，每个有限维对象都与某个 $\F^n$ 同构。
\end{remark}

\subsection{反范畴与函子}

为了统一处理“逆变”现象，先引入反范畴。

\begin{definition}
设 $\mathcal C$ 为范畴。其\emph{反范畴} $\mathcal C^{\mathrm{op}}$ 定义如下：对象与 $\mathcal C$ 相同，而
\[
  \Hom_{\mathcal C^{\mathrm{op}}}(X,Y)
  =
  \Hom_{\mathcal C}(Y,X).
\]
也就是说，$\mathcal C^{\mathrm{op}}$ 中的态射方向与 $\mathcal C$ 中完全相反，复合按相反顺序定义。
\end{definition}

\begin{definition}
设 $\mathcal C,\mathcal D$ 为范畴。一个（协变）\emph{函子}
\[
  F\colon \mathcal C\longrightarrow \mathcal D
\]
由以下数据组成：
\begin{enumerate}[label=(\arabic*), leftmargin=2.5em]
  \item 对每个对象 $X\in \mathcal C$，指定一个对象 $F(X)\in \mathcal D$；
  \item 对每个态射 $f\colon X\to Y$，指定一个态射
  \[
    F(f)\colon F(X)\to F(Y)
  \]
\end{enumerate}
满足：
\begin{enumerate}[label=(\roman*), leftmargin=2.5em]
  \item $F(g\circ f)=F(g)\circ F(f)$；
  \item $F(\Id_X)=\Id_{F(X)}$。
\end{enumerate}

一个\emph{逆变函子} $F\colon \mathcal C\to \mathcal D$ 指的是一个协变函子
\[
  F\colon \mathcal C^{\mathrm{op}}\longrightarrow \mathcal D.
\]
等价地说，它对态射 $f\colon X\to Y$ 赋予一个态射 $F(f)\colon F(Y)\to F(X)$，并满足
\[
  F(g\circ f)=F(f)\circ F(g),
  \qquad
  F(\Id_X)=\Id_{F(X)}.
\]
\end{definition}

\begin{example}[忘却函子]
从 $\mathbf{Grp}$ 到 $\mathbf{Set}$ 有一个忘却函子
\[
  U\colon \mathbf{Grp}\to \mathbf{Set},
\]
它把一个群送到其底层集合，把群同态看成集合映射。类似地，有
\[
  \mathbf{Ring}\to \mathbf{Set},
  \qquad
  R\text{-}\mathbf{Mod}\to \mathbf{Ab},
  \qquad
  \mathbf{Top}\to \mathbf{Set}
\]
等忘却函子。它们“忘记了一部分结构，但保留了态射”。
\end{example}

\begin{example}[自由函子]
忘却函子通常有左伴随，这里的典型例子是自由函子。例如从集合范畴到群范畴有自由群函子
\[
  F\colon \mathbf{Set}\to \mathbf{Grp},
\]
其中 $F(S)$ 是由集合 $S$ 生成的自由群。它满足泛性质：任意集合映射 $S\to U(G)$ 唯一延拓为群同态 $F(S)\to G$。类似地，对环 $R$，从 $\mathbf{Set}$ 到 $R\text{-}\mathbf{Mod}$ 有自由 $R$-模函子 $S\mapsto R^{(S)}$。
\end{example}

\begin{example}[Hom 函子]
设 $\mathcal C$ 为范畴，固定对象 $A\in \mathcal C$。则有两个与 $A$ 有关的 Hom 函子：
\begin{align*}
  \Hom_{\mathcal C}(A,-)&\colon \mathcal C\to \mathbf{Set},\\
  \Hom_{\mathcal C}(-,A)&\colon \mathcal C^{\mathrm{op}}\to \mathbf{Set}.
\end{align*}
前者把对象 $X$ 送到集合 $\Hom_{\mathcal C}(A,X)$；若 $f\colon X\to Y$，则
\[
  \Hom_{\mathcal C}(A,f)\colon \Hom_{\mathcal C}(A,X)\to \Hom_{\mathcal C}(A,Y),
  \qquad u\mapsto f\circ u.
\]
后者是逆变的：若 $f\colon X\to Y$，则
\[
  \Hom_{\mathcal C}(f,A)\colon \Hom_{\mathcal C}(Y,A)\to \Hom_{\mathcal C}(X,A),
  \qquad v\mapsto v\circ f.
\]
在加法范畴（如模范畴）中，Hom 甚至取值于阿贝尔群范畴。
\end{example}

\begin{example}[对偶函子]
在 $\mathbf{Vect}_{\F}$ 中，每个向量空间 $V$ 的对偶空间定义为
\[
  V^{\vee}=\Hom_{\F}(V,\F).
\]
若 $f\colon V\to W$ 为线性映射，则其对偶映射
\[
  f^{\vee}\colon W^{\vee}\to V^{\vee},
  \qquad \lambda\mapsto \lambda\circ f
\]
方向反转，因此 $V\mapsto V^{\vee}$ 给出一个逆变函子
\[
  (-)^{\vee}\colon \mathbf{Vect}_{\F}\to \mathbf{Vect}_{\F}.
\]
\end{example}

\begin{remark}
许多重要构造本质上都是函子：取商、取子对象、张量积、取不变量、取上同调群，等等。范畴语言的价值在于：一旦确认某个构造是函子，我们就知道它自动与态射复合相容。
\end{remark}

\section{自然变换}

若函子是“在范畴之间传送对象和态射的规则”，那么自然变换就是“在两个函子之间进行比较的规则”。

\begin{definition}
设 $F,G\colon \mathcal C\to \mathcal D$ 为两个协变函子。一个从 $F$ 到 $G$ 的\emph{自然变换}
\[
  \eta\colon F\Rightarrow G
\]
是指：对每个对象 $X\in \mathcal C$，给定一个态射
\[
  \eta_X\colon F(X)\to G(X),
\]
并且对任意态射 $f\colon X\to Y$，下图交换：
\[
\begin{tikzcd}
F(X) \arrow[r,"\eta_X"] \arrow[d,"F(f)"'] & G(X) \arrow[d,"G(f)"] \\
F(Y) \arrow[r,"\eta_Y"'] & G(Y)
\end{tikzcd}
\]
即
\[
  G(f)\circ \eta_X = \eta_Y\circ F(f).
\]
若每个 $\eta_X$ 都是同构，则称 $\eta$ 为\emph{自然同构}，并记 $F\cong G$。
\end{definition}

\begin{remark}
自然性的核心不是“存在某些同构”，而是“这些同构对所有态射同时相容”。很多看似显然的逐对象构造，若不满足这个交换图，就不是自然变换。
\end{remark}

\begin{example}[双对偶自然变换]
在 $\mathbf{Vect}_{\F}$ 中有从恒等函子到双对偶函子的自然变换
\[
  \eta\colon \Id \Rightarrow (-)^{\vee\vee}.
\]
其在对象 $V$ 上的分量定义为
\[
  \eta_V\colon V\to V^{\vee\vee},
  \qquad
  \eta_V(v)(\lambda)=\lambda(v)
  \quad (v\in V,\ \lambda\in V^{\vee}).
\]
\end{example}

\begin{proposition}
上面的 $\eta\colon \Id \Rightarrow (-)^{\vee\vee}$ 的确是一个自然变换。若 $V$ 有限维，则 $\eta_V$ 是同构；一般地，$\eta_V$ 总是单射。
\end{proposition}

\begin{proof}
先验证自然性。设 $f\colon V\to W$ 为线性映射。需要证明
\[
  f^{\vee\vee}\circ \eta_V = \eta_W\circ f.
\]
对任意 $v\in V$ 与任意 $\mu\in W^{\vee}$，有
\[
  \bigl((f^{\vee\vee}\circ \eta_V)(v)\bigr)(\mu)
  = \eta_V(v)(f^{\vee}(\mu))
  = f^{\vee}(\mu)(v)
  = \mu(f(v)).
\]
另一方面，
\[
  \bigl((\eta_W\circ f)(v)\bigr)(\mu)
  = \eta_W(f(v))(\mu)
  = \mu(f(v)).
\]
故二者相等，自然性成立。

再证单射。若 $v\in V$ 满足 $\eta_V(v)=0$，则对任意 $\lambda\in V^{\vee}$ 有 $\lambda(v)=0$。若 $v\neq 0$，可取一个基并把 $v$ 延拓为基的一员，再定义某个线性泛函使得它在 $v$ 上取值 $1$，矛盾。因此 $v=0$，故 $\eta_V$ 单射。

若 $V$ 有限维，设 $\dim_{\F}V=n$。则
\[
  \dim_{\F}V^{\vee}=n,
  \qquad
  \dim_{\F}V^{\vee\vee}=n.
\]
既然 $\eta_V$ 是有限维向量空间之间的单射，它必为同构。
\end{proof}

\begin{example}
设 $R$ 为环，$M$ 为固定左 $R$-模。则
\[
  \End_R(M)=\Hom_R(M,M)
\]
是一个环，而每个 $R$-模同构 $f\colon M\to N$ 都诱导一个环同构
\[
  \End_R(M)\to \End_R(N),
  \qquad
  u\mapsto f\circ u\circ f^{-1}.
\]
这说明很多代数对象其实是由函子性控制的；在更高层面上，这种“由态射诱导结构变化”的规律往往由自然变换统一表达。
\end{example}

\subsection{函子范畴}

\begin{definition}
设 $\mathcal C,\mathcal D$ 为范畴。\emph{函子范畴} $[\mathcal C,\mathcal D]$ 定义为：
\begin{itemize}[leftmargin=2em]
  \item 对象是从 $\mathcal C$ 到 $\mathcal D$ 的函子；
  \item 态射是这些函子之间的自然变换；
  \item 两个自然变换的复合按分量逐点定义：若
  \[
    \eta\colon F\Rightarrow G,
    \qquad
    \mu\colon G\Rightarrow H,
  \]
  则
  \[
    (\mu\circ \eta)_X \coloneqq \mu_X\circ \eta_X.
  \]
\end{itemize}
\end{definition}

\begin{proposition}
上述定义确实给出了一个范畴 $[\mathcal C,\mathcal D]$。
\end{proposition}

\begin{proof}
要验证的是复合满足结合律并存在恒等态射。设
\[
  \eta\colon F\Rightarrow G,
  \quad
  \mu\colon G\Rightarrow H,
  \quad
  \nu\colon H\Rightarrow K.
\]
对每个对象 $X\in \mathcal C$，有
\[
  ((\nu\circ \mu)\circ \eta)_X
  = (\nu_X\circ \mu_X)\circ \eta_X
  = \nu_X\circ (\mu_X\circ \eta_X)
  = (\nu\circ (\mu\circ \eta))_X,
\]
故结合律成立。对每个函子 $F$，取分量为恒等态射的自然变换
\[
  (\Id_F)_X\coloneqq \Id_{F(X)}.
\]
显然它满足自然性，而且对任意 $\eta\colon F\Rightarrow G$ 都有
\[
  \eta\circ \Id_F = \eta,
  \qquad
  \Id_G\circ \eta = \eta.
\]
因此 $[\mathcal C,\mathcal D]$ 构成一个范畴。
\end{proof}

\begin{remark}
当 $\mathcal D=\mathbf{Set}$ 时，函子范畴 $[\mathcal C^{\mathrm{op}},\mathbf{Set}]$ 的对象常称为 $\mathcal C$ 上的\emph{预层}。在代数几何与模形式理论中，层与预层语言极其重要。
\end{remark}

\subsection{Yoneda 观点的一瞥}

本节只点出一个非常重要的思想：对象常常可以通过它与所有其他对象之间的 Hom 集来刻画。

\begin{proposition}
设 $\mathcal C$ 为任意范畴，$A,B\in \mathcal C$。若存在自然同构
\[
  \Hom_{\mathcal C}(A,-)\cong \Hom_{\mathcal C}(B,-),
\]
则 $A\cong B$。
\end{proposition}

\begin{proof}
令
\[
  \Phi\colon \Hom_{\mathcal C}(A,-)\Rightarrow \Hom_{\mathcal C}(B,-)
\]
为给定的自然同构。取对象 $A$，则
\[
  \Phi_A\colon \Hom_{\mathcal C}(A,A)\xrightarrow{\sim} \Hom_{\mathcal C}(B,A).
\]
设
\[
  u\coloneqq \Phi_A(\Id_A)\colon B\to A.
\]
同理，利用逆自然同构 $\Phi^{-1}$ 在对象 $B$ 上的分量，得到
\[
  v\coloneqq \Phi^{-1}_B(\Id_B)\colon A\to B.
\]
我们来证 $u$ 与 $v$ 互为逆。由 $\Phi$ 的自然性，对态射 $v\colon A\to B$，交换图
\[
\begin{tikzcd}
\Hom(A,A) \arrow[r,"\Phi_A"] \arrow[d,"\Hom(A,v)"'] & \Hom(B,A) \arrow[d,"\Hom(B,v)"] \\
\Hom(A,B) \arrow[r,"\Phi_B"'] & \Hom(B,B)
\end{tikzcd}
\]
告诉我们
\[
  \Phi_B(v)=v\circ u.
\]
但由 $v=\Phi_B^{-1}(\Id_B)$，得 $\Phi_B(v)=\Id_B$，故 $v\circ u=\Id_B$。类似地对 $u\colon B\to A$ 应用 $\Phi^{-1}$ 的自然性，可得 $u\circ v=\Id_A$。所以 $A\cong B$。
\end{proof}

\begin{remark}
上述命题是 Yoneda 引理的一个简单推论。完整的 Yoneda 引理放在习题中引导证明。它说明：一个对象由其 Hom 函子完全决定。这是范畴论最深刻的基本原则之一。
\end{remark}

\section{正合序列与正合函子}

本节主要在模范畴 $R\text{-}\mathbf{Mod}$ 中讨论。对于阿贝尔群、向量空间、层以及更一般的阿贝尔范畴，也有同样的定义与结论。

\subsection{正合序列}

\begin{definition}
设
\[
  A \xrightarrow{f} B \xrightarrow{g} C
\]
为 $R$-模同态序列。若
\[
  \Ima(f)=\Ker(g),
\]
则称该序列在 $B$ 处\emph{正合}。

一个长序列
\[
  \cdots \to A_{i-1}\xrightarrow{d_{i-1}}A_i\xrightarrow{d_i}A_{i+1}\to \cdots
\]
若在每一项 $A_i$ 处都正合，则称其为\emph{正合序列}。
\end{definition}

\begin{definition}
若序列
\[
  0\longrightarrow A\xrightarrow{f} B\xrightarrow{g} C\longrightarrow 0
\]
正合，则称其为\emph{短正合序列}。
\end{definition}

\begin{proposition}
对于模同态序列
\[
  0\longrightarrow A\xrightarrow{f} B\xrightarrow{g} C\longrightarrow 0,
\]
其正合性等价于以下三条同时成立：
\begin{enumerate}[label=(\arabic*), leftmargin=2.5em]
  \item $f$ 是单射；
  \item $g$ 是满射；
  \item $\Ima(f)=\Ker(g)$。
\end{enumerate}
\end{proposition}

\begin{proof}
由定义，序列在 $A$ 处正合意味着
\[
  \Ker(f)=\Ima(0\to A)=0,
\]
故 $f$ 单射；在 $C$ 处正合意味着
\[
  \Ima(g)=\Ker(C\to 0)=C,
\]
故 $g$ 满射；在 $B$ 处正合正是 $\Ima(f)=\Ker(g)$。反过来，若这三条成立，则它分别保证序列在 $A,B,C$ 处正合，因此是短正合序列。
\end{proof}

\begin{example}
设 $R$ 为环，$I\subseteq R$ 为左理想。则
\[
  0\longrightarrow I\longrightarrow R\longrightarrow R/I\longrightarrow 0
\]
是短正合序列。
\end{example}

\begin{example}
设 $V$ 为向量空间，$W\subseteq V$ 为子空间。则
\[
  0\longrightarrow W\longrightarrow V\longrightarrow V/W\longrightarrow 0
\]
是 $\mathbf{Vect}_{\F}$ 中的短正合序列。
\end{example}

\subsection{左正合与右正合函子}

\begin{definition}
设 $F\colon R\text{-}\mathbf{Mod}\to \mathbf{Ab}$ 为函子。
\begin{enumerate}[label=(\arabic*), leftmargin=2.5em]
  \item 若对每个短正合序列
  \[
    0\to A\to B\to C\to 0
  \]
  ，施行 $F$ 后得到正合序列
  \[
    0\to F(A)\to F(B)\to F(C),
  \]
  则称 $F$ 为\emph{左正合函子}；
  \item 若施行 $F$ 后总得到正合序列
  \[
    F(A)\to F(B)\to F(C)\to 0,
  \]
  则称 $F$ 为\emph{右正合函子}；
  \item 若施行 $F$ 后总得到短正合序列
  \[
    0\to F(A)\to F(B)\to F(C)\to 0,
  \]
  则称 $F$ 为\emph{正合函子}。
\end{enumerate}
\end{definition}

\begin{remark}
严格地说，正合性最自然的背景是阿贝尔范畴以及加性函子。对本书而言，模范畴已经足够涵盖主要例子。
\end{remark}

\begin{proposition}[Hom 函子的左正合性]
设 $R$ 为环，$M$ 为左 $R$-模。则逆变函子
\[
  \Hom_R(-,M)\colon (R\text{-}\mathbf{Mod})^{\mathrm{op}}\to \mathbf{Ab}
\]
是左正合的。更明确地说，若
\[
  0\longrightarrow A\xrightarrow{u} B\xrightarrow{v} C\longrightarrow 0
\]
是左 $R$-模的短正合序列，则有正合序列
\[
  0\longrightarrow \Hom_R(C,M)
  \xrightarrow{v^{\ast}} \Hom_R(B,M)
  \xrightarrow{u^{\ast}} \Hom_R(A,M),
\]
其中
\[
  v^{\ast}(\varphi)=\varphi\circ v,
  \qquad
  u^{\ast}(\psi)=\psi\circ u.
\]
\end{proposition}

\begin{proof}
先证 $v^{\ast}$ 单射。设 $\varphi\in \Hom_R(C,M)$ 且 $v^{\ast}(\varphi)=0$，即 $\varphi\circ v=0$。由于 $v$ 满射，对任意 $c\in C$，可取 $b\in B$ 使 $v(b)=c$，于是
\[
  \varphi(c)=\varphi(v(b))=0.
\]
故 $\varphi=0$。

再证
\[
  \Ima(v^{\ast})=\Ker(u^{\ast}).
\]
若 $\varphi\in \Hom_R(C,M)$，则
\[
  u^{\ast}(v^{\ast}(\varphi))=(\varphi\circ v)\circ u=
  \varphi\circ (v\circ u)=0,
\]
所以 $\Ima(v^{\ast})\subseteq \Ker(u^{\ast})$。

反过来，取 $\psi\in \Hom_R(B,M)$，满足 $u^{\ast}(\psi)=0$，即 $\psi\circ u=0$。这表示 $\psi$ 在 $\Ima(u)=\Ker(v)$ 上为零。于是可以定义映射
\[
  \varphi\colon C\to M,
  \qquad
  \varphi(v(b))\coloneqq \psi(b).
\]
需证其良定义。若 $v(b_1)=v(b_2)$，则 $b_1-b_2\in \Ker(v)=\Ima(u)$，故存在 $a\in A$ 使 $b_1-b_2=u(a)$。于是
\[
  \psi(b_1)-\psi(b_2)=\psi(u(a))=(\psi\circ u)(a)=0.
\]
故 $\varphi$ 良定义。它显然是 $R$-模同态，且满足 $\psi=\varphi\circ v$，即 $\psi=v^{\ast}(\varphi)$。因此 $\Ker(u^{\ast})\subseteq \Ima(v^{\ast})$，正合性得证。
\end{proof}

\begin{remark}
一般来说，$\Hom_R(-,M)$ 并不右正合；其右导出函子正是 $\mathrm{Ext}_R^i(-,M)$。另一方面，协变函子 $\Hom_R(M,-)$ 也是左正合的，其右导出函子同样给出 Ext 群。
\end{remark}

\begin{proposition}[张量积的右正合性]
设 $R$ 为环，$N$ 为右 $R$-模。则函子
\[
  N\otimes_R -\colon R\text{-}\mathbf{Mod}\to \mathbf{Ab}
\]
是右正合的。也就是说，若
\[
  A\xrightarrow{u} B\xrightarrow{v} C\longrightarrow 0
\]
正合，则
\[
  N\otimes_R A\xrightarrow{1\otimes u} N\otimes_R B\xrightarrow{1\otimes v} N\otimes_R C\longrightarrow 0
\]
也正合。
\end{proposition}

\begin{proof}
先证 $(1\otimes v)$ 满射。任取 $N\otimes_R C$ 中的纯张量 $n\otimes c$。由于 $v$ 满射，可取 $b\in B$ 使 $v(b)=c$，于是
\[
  (1\otimes v)(n\otimes b)=n\otimes v(b)=n\otimes c.
\]
而任意张量都是纯张量的有限和，所以 $1\otimes v$ 满射。

再证
\[
  \Ima(1\otimes u)=\Ker(1\otimes v).
\]
由 $v\circ u=0$ 可知
\[
  (1\otimes v)\circ (1\otimes u)=1\otimes (v\circ u)=0,
\]
故 $\Ima(1\otimes u)\subseteq \Ker(1\otimes v)$。

为证反向包含，考虑商模 $B/u(A)$。由原序列正合知 $\Ker(v)=u(A)$，故由模的第一同构定理，有自然同构
\[
  \overline v\colon B/u(A)\xrightarrow{\sim} C.
\]
另一方面，张量积与余核相容：由张量积的泛性质可得自然同构
\[
  (N\otimes_R B)/\Ima(1\otimes u)\xrightarrow{\sim} N\otimes_R (B/u(A)).
\]
将其与 $\overline v$ 诱导的同构复合，得到
\[
  (N\otimes_R B)/\Ima(1\otimes u)\xrightarrow{\sim} N\otimes_R C.
\]
而这个同构正是由 $1\otimes v$ 诱导的，因此其核恰为 $\Ima(1\otimes u)$。即
\[
  \Ker(1\otimes v)=\Ima(1\otimes u).
\]
故序列右正合。
\end{proof}

\begin{remark}
张量积一般不左正合。若 $N$ 是平坦模，则 $N\otimes_R -$ 才是正合函子。平坦性在代数几何与模形式的整系数理论中经常出现。
\end{remark}

\subsection{长正合序列与蛇引理}

导出函子的基本作用，是把“左正合但不完全正合”或“右正合但不完全正合”的缺陷测量出来。例如，若 $F$ 是左正合函子，则其右导出函子 $R^iF$ 往往把短正合序列送成长正合序列：
\[
  0\to F(A)\to F(B)\to F(C)\to R^1F(A)\to R^1F(B)\to R^1F(C)\to \cdots.
\]
对 $F=\Hom_R(-,M)$，这条长正合序列给出 Ext 群；对层的全局截面函子 $\Gamma(X,-)$，其右导出函子给出层上同调群 $H^i(X,-)$。

\begin{theorem}[蛇引理]
设有交换图
\[
\begin{tikzcd}
0 \arrow[r] & A' \arrow[r] \arrow[d,"\alpha"] & B' \arrow[r] \arrow[d,"\beta"] & C' \arrow[r] \arrow[d,"\gamma"] & 0 \\
0 \arrow[r] & A \arrow[r] & B \arrow[r] & C \arrow[r] & 0
\end{tikzcd}
\]
其中两行都是短正合序列。则存在自然构造的连接同态
\[
  \delta\colon \Ker(\gamma)\longrightarrow \coker(\alpha),
\]
使得序列
\[
  0\to \Ker(\alpha)\to \Ker(\beta)\to \Ker(\gamma)
  \xrightarrow{\delta}
  \coker(\alpha)\to \coker(\beta)\to \coker(\gamma)\to 0
\]
正合。
\end{theorem}

\begin{remark}
蛇引理是导出长正合序列的标准工具之一。它把核与余核联系起来，许多同调代数中的“连接同态”都源于此。完整证明需要细致追踪图中的元素；本附录只陈述结论，相关练习会帮助读者熟悉其用法。
\end{remark}

\section{为什么模形式理论需要范畴语言}

若只把模形式看作上半平面上的某些解析函数，那么范畴语言似乎并非必需；但一旦进入现代观点，范畴语言就几乎无处不在。

\subsection*{(1) 模形式是层的整体截面}

设 $X$ 是适当的模曲线，$\omega$ 是其上的 Hodge 线丛。权 $k$ 的模形式空间可以统一写成
\[
  M_k(\Gamma) = H^0(X,\omega^{\otimes k}),
\]
尖点形式则对应于加入尖点条件后的截面子空间。这里 $H^0(X,-)$ 不是简单的“取函数”，而是一个从层范畴到阿贝尔群范畴的函子。于是模形式本身应被看作某个层对象的全局截面，而不是孤立的函数集合。

\subsection*{(2) Hecke 算子来自对应的函子性}

Hecke 算子并不是凭空写下来的线性变换。对模曲线之间的 Hecke 对应，通常存在形如
\[
  X \xleftarrow{\pi_1} Y \xrightarrow{\pi_2} X
\]
的两个有限态射。它们对层及其上同调诱导拉回与推前：
\[
  \pi_2^{\ast},\qquad (\pi_1)_{\ast}.
\]
Hecke 算子正是这些态射经复合得到的线性算子。没有函子与自然变换的语言，就很难清楚说明这些构造为何彼此兼容、为何与基变换相容、为何在不同实现之间给出同一个算子。

\subsection*{(3) 上同调与导出函子刻画更深层结构}

Eichler--Shimura 理论把权 $2$ 以上的模形式与模曲线上的上同调联系起来。粗略地说，尖点形式空间可以嵌入某个一阶上同调群中，而 Hecke 算子在两边兼容。这种叙述本质上属于导出函子与长正合序列的世界：
\[
  H^i(X,\mathscr F)=R^i\Gamma(X,\mathscr F).
\]
若没有“函子、导出函子、自然性”的语言，我们很难精确表述 Eichler--Shimura 同构、比较同构以及各种 $\ell$-进上同调实现之间的关系。

\subsection*{(4) 现代数论中的统一框架}

现代模形式理论不仅研究复解析模形式，还研究整系数模形式、$p$-进模形式、代数模形式、Galois 表示与自守表示之间的关系。这些对象分布在不同范畴中：
\begin{itemize}[leftmargin=2em]
  \item 模曲线上的层与向量丛；
  \item Hecke 模；
  \item Galois 表示的范畴；
  \item 带有额外结构的模与复形；
  \item 导出范畴中的对象。
\end{itemize}
范畴语言提供了一个共同框架，使“同一数学对象的不同实现”能够通过函子加以比较。

\begin{remark}
对初学者而言，范畴语言的最低要求并不是立刻掌握导出范畴，而是先熟悉下面三个习惯：
\begin{enumerate}[label=(\arabic*), leftmargin=2.5em]
  \item 把研究对象放入一个合适的范畴；
  \item 识别常见构造是否给出函子；
  \item 在比较两个构造时，优先寻找自然变换而非逐对象的偶然同构。
\end{enumerate}
一旦形成这种思维方式，后续阅读模形式、代数几何与同调代数文献会容易得多。
\end{remark}

\section*{习题}

下列习题按难度分为三级：\basic\\ 为基础题，\medium\\ 为进阶题，\hard\\ 为挑战题。

\subsection*{基础题}

\begin{enumerate}[label=\textbf{习题A.\arabic*}, leftmargin=3.5em]
  \item \basic\ 设 $\mathcal C$ 的对象只有两个 $X,Y$，且态射只有 $\Id_X,\Id_Y$ 与一个态射 $f\colon X\to Y$。定义所有可能的复合，并验证这确实给出一个范畴。

  \item \basic\ 证明：在任意范畴中，每个对象的恒等态射是唯一的。

  \item \basic\ 设 $f\colon X\to Y$ 在某范畴中可逆。证明其逆态射唯一。

  \item \basic\ 在下列构造中，判断它是否给出函子；若是，指出它是协变还是逆变：
  \begin{enumerate}[label=(\alph*), leftmargin=2.5em]
    \item $\mathbf{Vect}_{\F}\to \mathbf{Set}$，把向量空间送到底层集合；
    \item $\mathbf{Vect}_{\F}\to \mathbf{Vect}_{\F}$，把 $V$ 送到 $V^{\vee}$；
    \item $R\text{-}\mathbf{Mod}\to \mathbf{Ab}$，把 $M$ 送到 $M/2M$；
    \item $\mathbf{Top}\to \mathbf{Grp}$，把拓扑空间送到其基本群（只需忽略基点问题，讨论思路即可）。
  \end{enumerate}

  \item \basic\ 设 $A$ 为范畴 $\mathcal C$ 的固定对象。写出 Hom 函子 $\Hom_{\mathcal C}(A,-)$ 与 $\Hom_{\mathcal C}(-,A)$ 在态射上的具体作用，并验证函子公理。

  \item \basic\ 在 $\mathbf{Set}$ 中，设 $F$ 为“取幂集”函子 $X\mapsto \mathcal P(X)$，$G$ 为恒等函子。证明不存在自然变换 $\mathcal P\Rightarrow \Id$，但存在自然变换 $\Id\Rightarrow \mathcal P$。

  \item \basic\ 设 $V,W$ 为有限维向量空间。证明自然映射
  \[
    \Hom_{\F}(V,W)\longrightarrow \Hom_{\F}(W^{\vee},V^{\vee}),
    \qquad f\mapsto f^{\vee}
  \]
  是向量空间同构。

  \item \basic\ 设 $0\to A\xrightarrow{u}B\xrightarrow{v}C\to 0$ 为短正合序列。证明若该序列分裂，即存在 $s\colon C\to B$ 满足 $v\circ s=\Id_C$，则 $B\cong A\oplus C$。

  \item \basic\ 在 $\Z$-模范畴中，验证
  \[
    0\to \Z \xrightarrow{\times 2} \Z \to \Z/2\Z \to 0
  \]
  是短正合序列。

  \item \basic\ 设 $M$ 为固定 $R$-模。证明赋值 $N\mapsto \Hom_R(M,N)$ 连同态射上的后复合，给出一个协变函子 $R\text{-}\mathbf{Mod}\to \mathbf{Ab}$。
\end{enumerate}

\subsection*{进阶题}

\begin{enumerate}[resume, label=\textbf{习题A.\arabic*}, leftmargin=3.5em]
  \item \medium\（Yoneda 引理，分步证明）设 $\mathcal C$ 为范畴，$F\colon \mathcal C\to \mathbf{Set}$ 为协变函子，$A\in \mathcal C$。证明映射
  \[
    \mathrm{Nat}(\Hom_{\mathcal C}(A,-),F)\longrightarrow F(A),
    \qquad \eta\mapsto \eta_A(\Id_A)
  \]
  是双射。提示：
  \begin{enumerate}[label=(\alph*), leftmargin=2.5em]
    \item 对每个 $x\in F(A)$，构造自然变换 $\eta^x$；
    \item 验证 $\eta^x_X(f)=F(f)(x)$；
    \item 证明上述构造互为逆。
  \end{enumerate}

  \item \medium\ 利用上一题证明：若 $\Hom_{\mathcal C}(A,-)$ 与 $\Hom_{\mathcal C}(B,-)$ 自然同构，则 $A\cong B$。

  \item \medium\ 设 $0\to A'\to A\to A''\to 0$ 与 $0\to B'\to B\to B''\to 0$ 为两条短正合序列，且有交换图把前者映到后者。证明若左、右两端的竖箭头同构，则中间竖箭头也同构（五引理的短形式）。

  \item \medium\ 对短正合序列
  \[
    0\to \Z \xrightarrow{\times n} \Z \to \Z/n\Z \to 0
  \]
  施行函子 $\Hom_{\Z}(-,\Z)$，计算得到的序列，并说明它通常不是右正合的。

  \item \medium\ 设 $R$ 为环，$N$ 为右 $R$-模。补全正文中张量积右正合性的证明：直接由张量积的泛性质构造同构
  \[
    (N\otimes_R B)/\Ima(1\otimes u)\cong N\otimes_R(B/u(A)).
  \]

  \item \medium\ 设 $R=\Z$，将短正合序列
  \[
    0\to \Z \xrightarrow{\times 2} \Z \to \Z/2\Z \to 0
  \]
  与函子 $\Z/3\Z\otimes_{\Z}-$ 作用，计算所得序列，并据此验证右正合性。

  \item \medium\ 设 $R$ 为环，$M$ 为左 $R$-模。证明：若 $M$ 是投射模，则函子 $\Hom_R(M,-)$ 为正合函子。
\end{enumerate}

\subsection*{挑战题}

\begin{enumerate}[resume, label=\textbf{习题A.\arabic*}, leftmargin=3.5em]
  \item \hard\ 设 $U\colon \mathbf{Grp}\to \mathbf{Set}$ 为忘却函子，$F\colon \mathbf{Set}\to \mathbf{Grp}$ 为自由群函子。证明存在自然双射
  \[
    \Hom_{\mathbf{Grp}}(F(S),G)\cong \Hom_{\mathbf{Set}}(S,U(G)),
  \]
  并说明这意味着 $F$ 是 $U$ 的左伴随。

  \item \hard\ 设 $R$ 为环，证明函子 $\Hom_R(R,-)\colon R\text{-}\mathbf{Mod}\to \mathbf{Ab}$ 可表表示，并写出表示对象。进一步说明：若 $P$ 是有限生成投射模，则函子 $\Hom_R(P,-)$ 在什么意义下“接近可表表示”。

  \item \hard\ 设 $m,n\ge 1$。利用短正合序列
  \[
    0\to \Z \xrightarrow{\times n} \Z \to \Z/n\Z \to 0
  \]
  和函子 $\Hom_{\Z}(-,\Z/m\Z)$，推导
  \[
    \mathrm{Ext}^1_{\Z}(\Z/n\Z,\Z/m\Z)\cong \Z/\gcd(m,n)\Z.
  \]
  提示：先写出长正合序列，再计算乘以 $n$ 在 $\Z/m\Z$ 上的像与核。
\end{enumerate}
